% ============================================================================
% CHAPITRE 4 — SECTION 4.4 : PERSPECTIVES DE RECHERCHE
% Fichier : partie4/perspectives.tex
% Version : 2.2 — alignée sur la formalisation Lean4 révisée
% ============================================================================

\section{PERSPECTIVES DE RECHERCHE}
\label{sec:perspectives}
% ============================================================================
\begin{sloppypar}
Le corpus th\'eorique de la cryptographie institutionnelle est maintenant
constitu\'e. Il r\'eunit des m\'etriques formalis\'ees (C1)~\cite{MINKA2026CLO}, une borne
d'impossibilit\'e \'etablie (C2)~\cite{MINKA2026CRO}, une extension quantique
d\'emontr\'ee (C3)~\cite{MINKA2026PQCRO}, des constructions de r\'ef\'erence aux propri\'et\'es
cryptographiques prouv\'ees (C4) (section~\ref{sec:sotp_construction}), un noyau Lean4 certifi\'e (C5)
(section~\ref{sec:lean4}) et un programme de validation sp\'ecifi\'e (C6)
(chapitre~\ref{chap:experimentation}). Que faire de ces fondations~?
Cette section y r\'epond en
tra\c{c}ant la feuille de route th\'eorique d'un sous-domaine dont les
premiers th\'eor\`emes viennent d'\^etre pos\'es, sans projeter de d\'eploiements.

La progression F$\to$B$\to$E$\to$S$\to$C$\to$X
(formaliser, borner, \'etendre, saturer, certifier, eXp\'erimenter)
qui a guid\'e cette th\`ese se prolonge au-del\`a de ces r\'esultats, comme le montrent les quatre axes qui structurent cette
prospective.

Chacun de ces quatre axes prolonge un aspect
distinct du cadre CLO-CRO-PQ-SOTP. Il s'agit des fondations computationnelles
post-quantiques port\'ees par le probl\`eme AIIP~\cite{MINKA2025AIIP} et le
CASH Toolkit (section~\ref{subsec:axe1_aiip_cash})~; la voie protocolaire
ouverte par Q2CSI et la s\'erie
ZK-NR~\cite{MINKA2025Q2CSI,MINKA2025ZKNR_Exploration,MINKA2025ZKNR_Design,MINKA2025ZKNR_UC}
(section~\ref{subsec:axe2_q2csi_zknr})~; les extensions th\'eoriques du
Trilemme CRO vers le calcul homomorphe, le multi-parties et le
post-quantique quantitatif (section~\ref{subsec:axe3_extensions})~; et
l'interaction encore largement inexploit\'ee entre cryptographie
institutionnelle et apprentissage automatique (section~\ref{subsec:axe4_ia}).
Un programme \`a horizons structur\'es conclut la section
(section~\ref{subsec:programme_recherche}).
\end{sloppypar}

\begin{boiteArbitre}
Le voyage de l'arbitre ne s'arr\^ete pas avec cette th\`ese. Les fondations
sont pos\'ees (m\'etriques, bornes, constructions, certification,
protocoles) mais le programme ouvert est plus vaste que ce qui vient
d'\^etre r\'esolu. Il reste \`a obtenir une borne dont le terme dominant serait en $O(\lambda^{-2})$,
\`a \'etablir l'optimalit\'e des constructions SOTP, \`a calibrer empiriquement les param\`etres
quantiques~$\Delta_q$ et~$\eta_q$ ainsi qu'\`a \'etendre les constructions
au calcul homomorphe et multi-parties. L'arbitre de demain
aura besoin de protocoles que la cryptographie institutionnelle ne sait pas
encore concevoir. Ce chapitre en trace la feuille de route.
\end{boiteArbitre}

% -----------------------------------------------------------------------
\subsection{Maturit\'e du sous-domaine}
\label{subsec:maturite_sous_domaine}
% -----------------------------------------------------------------------
\begin{sloppypar}
Un sous-domaine scientifique est constitu\'e lorsqu'il dispose d'un formalisme
propre, d'une impossibilit\'e fondamentale, de constructions et d'une
certification m\'ecanique~\cite{Wiedijk2007,Gonthier08}. Le cadre
CLO-CRO-CRO$^q$-SOTP-Lean4 r\'eunit ces \'el\'ements, sous deux r\'eserves. D'une part,
l'optimalit\'e des constructions est conjecturale~; d'autre part, les param\`etres
institutionnels restent \`a mesurer. Il s'agit d'un programme de recherche au sens
de Lakatos~\cite{Lakatos1978}, dont le noyau est fix\'e et la ceinture
protectrice explicite.

La question devient \guillemotleft~quelles directions ouvertes en
d\'efiniront la croissance et quelles mesures en confirmeront ou en
infirmeront les hypoth\`eses~?~\guillemotright.

L'\'etat des contributions \`a la date de soutenance figure au
tableau~\ref{tab:contributions_these}.

Trois classes de probl\`emes ouverts structurent la fronti\`ere du
sous-domaine. La premi\`ere concerne les fondations computationnelles, c'est-\`a-dire la n\'ecessit\'e
d'asseoir les instanciations pratiques sur des hypoth\`eses de duret\'e
nouvelles, r\'esistantes aux adversaires quantiques~\cite{Bernstein2009,Chen2016}
et analys\'ees par la communaut\'e IACR. La deuxi\`eme concerne la
g\'en\'eralisation des bornes. Dans ses formes classique et
quantique, le Trilemme CRO couvre les preuves non interactives~\cite{Groth16}~; le passage
aux protocoles \`a rounds multiples, au calcul multi-parties~\cite{Evans2018}
et aux primitives homomorphes~\cite{Gentry09} reste entier. La troisi\`eme
est l'ouverture d'un territoire in\'edit, car l'interaction entre interpr\`etes
institutionnels artificiels et protocoles SOTP modifie structurellement les
hypoth\`eses du mod\`ele.
\end{sloppypar}

% -----------------------------------------------------------------------
\subsection{Axe I. Fondations computationnelles~: AIIP et le CASH Toolkit}
\label{subsec:axe1_aiip_cash}
% -----------------------------------------------------------------------
\begin{sloppypar}

\subsubsection{Le probl\`eme AIIP~: hypoth\`ese de duret\'e post-quantique}
\label{subsubsec:probleme_aiip}

Toute discipline cryptographique repose sur des hypoth\`eses de duret\'e dont
le statut est \'etabli par cryptanalyse
communautaire~\cite{Goldreich01,KatzLindell14}. La cryptographie
institutionnelle ne fait pas exception. En effet, si les Trilemmes CRO classiques et
quantiques sont prouv\'es au sens de la th\'eorie de l'information, leurs
instanciations pratiques requi\`erent des primitives r\'esistant aux
adversaires quantiques \`a temps polynomial~(QPT)~\cite{Boneh2011,Watrous2009}.

Le \textbf{probl\`eme AIIP} (\textit{Affine Iterated Inversion
Problem})~\cite{MINKA2025AIIP} propose une telle hypoth\`ese. \'Etant
donn\'e $f \in \mathbb{F}_q[x]$ de degr\'e $d \geq 2$, un param\`etre
d'it\'eration $n$ et une cible $y \in \mathbb{F}_q$, AIIP requiert de trouver
$x$ tel que $f^{(n)}(x) = y$, o\`u $f^{(n)}$ d\'esigne la $n$-i\`eme
composition it\'er\'ee. Sa duret\'e s'\'etablit par deux r\'eductions
ind\'ependantes~\cite{MINKA2025AIIP}, \`a savoir une r\'eduction au logarithme discret dans
le Jacobien d'une courbe hyperelliptique de genre exponentiellement
grand~(HCDLP)~\cite{Menezes2001,Enge2002} et une r\'eduction polynomiale \`a la
r\'esolution de syst\`emes MQ structur\'es~(MQ-NP)~\cite{Courtois2000,Ding2004}.
Ces r\'eductions et leur port\'ee restent \`a \'etablir de fa\c{c}on
ind\'ependante. En particulier, une objection doit \^etre lev\'ee avant de
consid\'erer AIIP comme une hypoth\`ese candidate. En effet, sauf argument contraire,
r\'esoudre $f^{(n)}(x) = y$ se ram\`ene \`a $n$ extractions successives de racines
d'un polyn\^ome de degr\'e $d$ sur $\mathbb{F}_q$, ce qui se fait en temps
polynomial en $n$, $d$ et $\log q$ (algorithme de Cantor--Zassenhaus), en
suivant les ant\'ec\'edents successifs. Un argument qui exclut cette attaque
(par exemple un contr\^ole de la ramification de l'arbre des ant\'ec\'edents)
est donc un pr\'ealable. La cryptanalyse (attaques de
Gr\"obner~\cite{Faugere02,Faugere99}, algorithmes de
point-count~\cite{Schoof85,Kedlaya2001}, variantes
quantiques~(Grover~\cite{Grover96}, MNFS quantique~\cite{Jeppesen2022}))
reste \`a conduire.

\subsubsection{Le CASH Toolkit~: primitives align\'ees sur les propri\'et\'es CLO}
\label{subsubsec:cash_toolkit}

Le \textbf{CASH Toolkit} (\textit{Chaotic Affine Secure Hash}) regroupe trois
primitives post-quantiques construites en correspondance directe avec les
dimensions du tripl\'e~$(\Conf, \Rel, \Opp)$. Leur conception est guid\'ee par
un principe inhabituel dans la cryptographie classique, selon lequel l'opposabilit\'e
future de la preuve conditionne d\`es la conception les choix de transparence
alg\'ebrique~\cite{MINKA2026CLO,Casey11}.

\paragraph{CEE. Couche confidentialit\'e ($\Conf$).}
La primitive \textbf{CEE} (\textit{Chaotic Entropic
Expansion})~\cite{MINKA2025CEE} construit une fonction \`a sens unique sur des
applications polynomiales chaotiques it\'er\'ees sur
$\mathbb{F}_q$~\cite{Menezes1996,Stinson2005}. Pour $N$ it\'erations d'un
polyn\^ome $f \in \mathcal{F}_d$, la min-entropie r\'esiduelle est au moins
$\log_2 q - N\log_2 d$ bits. La borne de distance statistique \`a l'uniforme
annonc\'ee dans~\cite{MINKA2025CEE}, $(q-1)(d^N-1)/(2\sqrt{q})$, d\'epasse $1$ d\`es
que $d^N > 2\sqrt{q}/(q-1) + 1$, c'est-\`a-dire pour tout $N$ utile. Elle est
alors vide, si bien qu'une borne non triviale reste \`a \'etablir. La s\'ecurit\'e
quantique est de l'ordre $O(2^{\lambda/2})$ pour un niveau classique
$\lambda$~\cite{NielsenChuang2010}. La construction a pour trait caract\'eristique que sa transparence
alg\'ebrique est d\'elib\'er\'ement pr\'eserv\'ee, permettant
\`a $\Conf$ de coexister avec $\Opp$ sans conflit structurel~\cite{MINKA2026CLO}.

\paragraph{AOW. Couche fiabilit\'e ($\Rel$).}
La primitive \textbf{AOW} (\textit{Affine One-Wayness})~\cite{MINKA2025AOW}
fournit une v\'erification temporelle par \'evaluation polynomiale it\'erative
sur corps finis. Sa s\'ecurit\'e est r\'eduite, dans~\cite{MINKA2025AOW}, \`a HCDLP
et \`a AIIP, dont le statut est discut\'e ci-dessus~; elle s'accompagne d'une int\'egration
STARK~\cite{STARKs18,BenSasson2018} \`a mise \`a l'\'echelle logarithmique
pour la v\'erification zero-knowledge. La tra\c{c}abilit\'e temporelle qu'AOW
produit, ancrage de s\'equences v\'erifiables sans autorit\'e de
confiance~\cite{Haber1991,Bayer1993}, alimente directement $\Opp$, puisque
l'institution peut \'etablir \emph{quand} une preuve a \'et\'e g\'en\'er\'ee.

\paragraph{SH. Couche opposabilit\'e ($\Opp$).}
La primitive \textbf{SH} (\textit{Semantic Holder})~\cite{MINKA2025SH}
op\'erationnalise l'extraction alg\'ebrique depuis les traces de polyn\^omes
it\'er\'es pour produire des interpr\'etations l\'egalement
opposables~\cite{Kerr04,Casey11}. Elle vise trois propri\'et\'es, \`a savoir la
pr\'eservation entropique, l'explicabilit\'e institutionnelle et la contestabilit\'e
l\'egale v\'erifiable par STARK. Le score agr\'eg\'e $\Omega$ et la valeur
$\Omega \geq 0{,}60$ avanc\'es dans~\cite{MINKA2025SH} ne sont pas d\'efinis dans le
cadre de cette th\`ese~; ils ne sont pas repris ici comme r\'esultats.

Plusieurs directions restent ouvertes. Peut-on prouver que l'assemblage
$(CEE \;\|\; AOW \;\|\; SH)$ en architecture Gold-Iron-Clay \emph{pr\'eserve}
les bornes CRO \`a un facteur constant pr\`es~? Peut-on d\'efinir un score
d'opposabilit\'e pour SH qui soit reli\'e \`a $\Opp$~? Enfin, les preuves actuelles
op\'erant dans le mod\`ele de l'oracle al\'eatoire~\cite{Bellare1993,Canetti2004},
des preuves en mod\`ele standard renforceraient substantiellement la fondation
th\'eorique du toolkit.
\end{sloppypar}

% -----------------------------------------------------------------------
\subsection{Axe II. Voie protocolaire~: Q2CSI, ZK-NR et l'instanciation de l'esprit SOTP}
\label{subsec:axe2_q2csi_zknr}
% -----------------------------------------------------------------------
\begin{sloppypar}

\subsubsection{Q2CSI comme traduction architecturale de l'esprit SOTP}
\label{subsubsec:q2csi}

La famille SOTP propose des constructions de r\'ef\'erence par r\'egime,
param\'etr\'ees par $H(\mathcal{C})$ (section~\ref{sec:sotp_construction}), dont
l'optimalit\'e reste conjecturale. La question qui
s'ouvre naturellement est celle de l'\emph{instanciation concr\`ete}. Comment
traduire cette philosophie en architectures d\'eployables dans des
environnements cryptographiques composables~\cite{Canetti01,Canetti2000}~?

Le m\'eta-protocole \textbf{Q2CSI} (\textit{Quantum Composable and Contextual
Security Infrastructure})~\cite{MINKA2025Q2CSI} apporte une r\'eponse
pr\'ecise. En d\'ecomposant les garanties de s\'ecurit\'e en trois couches
fonctionnellement isol\'ees (Iron~($\Rel$), Gold~($\Conf$),
Clay~($\Opp$)) Q2CSI traduit architecturalement le paradigme
Gold-Iron-Clay dans le mod\`ele de Canetti~\cite{Canetti01,Canetti2000}.
L'article~\cite{MINKA2025Q2CSI} r\'esume ses performances par un indice
$\Gamma_\text{CRO}$ qui n'est pas d\'efini dans le cadre de cette th\`ese~; les
valeurs chiffr\'ees qu'il en donne ne sont donc pas reprises ici. Relier
cet indice au produit $\Conf\cdot\Relrec\cdot\Opp$ est une question ouverte. Q2CSI incarne
l'esprit SOTP sous la forme d'un sch\'ema d'architecture, distinct d'une
construction abstraite. Ce sch\'ema repose sur la s\'eparation dialectique des couches, l'isolation stricte des
fuites contextuelles et la compatibilit\'e avec les adversaires
quantiques~\cite{Unruh2010,Unruh2015}.

\subsubsection{La s\'erie ZK-NR~: instanciation composable par zero-knowledge}
\label{subsubsec:zk_nr}

La s\'erie \textbf{ZK-NR} (\textit{Zero-Knowledge Non-Repudiation}) protocol
mat\'erialise la maturation progressive d'un protocole institutionnel depuis
l'exploration architecturale jusqu'\`a la s\'ecurit\'e UC
compl\`ete~\cite{Canetti01}. Elle constitue en elle-m\^eme un paradigme
m\'ethodologique, qui consiste \`a s\'eparer l'exploration, la formalisation du design et la
certification de s\'ecurit\'e en trois stades distincts, chacun s'appuyant
rigoureusement sur le pr\'ec\'edent.

ZK-NR~v1~\cite{MINKA2025ZKNR_Exploration} compose commitments
Merkle~\cite{Merkle1987}, preuves STARKs~\cite{STARKs18}, signatures BLS \`a
seuil~\cite{Boldyreva2003,Boneh2003} et authentification
Dilithium~\cite{Dilithium2022} en un protocole de non-r\'epudiation sans
divulgation~; il \'etablit la faisabilit\'e conceptuelle, les preuves formelles
\'etant r\'eserv\'ees aux versions ult\'erieures. ZK-NR~v2~\cite{MINKA2025ZKNR_Design}
construit l'architecture compl\`ete satisfaisant les propri\'et\'es Q2CSI et le
mod\`ele NIZK-E~\cite{Groth16}, avec s\'eparation structurelle entre preuves
contextuelles et explications born\'ees~; la compatibilit\'e UC est prouv\'ee
symboliquement. ZK-NR~v3~\cite{MINKA2025ZKNR_UC} \'etablit la preuve
compl\`ete de s\'ecurit\'e UC pour les trois couches, mod\'elis\'ees comme
fonctionnalit\'es id\'eales avec s\'emantique d'effacement~\cite{Canetti2016}.
Il prouve aussi l'indistinguabilit\'e r\'eel/id\'eal sous hypoth\`eses
post-quantiques~\cite{Unruh2015}.

Trois questions ouvertes d\'elimitent la suite. Le sch\'ema Q2CSI est-il
\emph{complet}, c'est-\`a-dire, tout protocole institutionnel qui s\'epare
les trois couches est-il UC-\'equivalent \`a une instanciation Q2CSI~? \`a l'image du th\'eor\`eme de compl\'etude de
Canetti-Cohen-Lindell~\cite{CCL15}. La r\'eduction formelle
SOTP-High~$\to$~ZK-NR~v3 est-elle exhibable, comblant l'\'ecart entre la
th\'eorie abstraite de SOTP~(section~\ref{sec:sotp_construction}) et les
protocoles composables d\'eployables~? La structure \`a couches de ZK-NR
s'\'etend-elle aux r\'eseaux de preuves institutionnelles multi-parties
impliquant $n > 2$ institutions, via des sch\'emas de partage de
secret~\cite{Shamir1979,Blakley1979} et des preuves zero-knowledge
agr\'eg\'ees~\cite{Bunz2018}~?
\end{sloppypar}

% -----------------------------------------------------------------------
\subsection{Axe III. Extensions th\'eoriques du Trilemme CRO}
\label{subsec:axe3_extensions}
% -----------------------------------------------------------------------
\begin{sloppypar}

\subsubsection{Borne forte sur le produit CRO~: vers la caract\'erisation compl\`ete}
\label{subsubsec:borne_forte}

Le Trilemme CRO classique~(Th\'eor\`eme~\ref{thm:trilemme_cro_conditionnel})
\'etablit la borne explicite
\begin{equation}
\Conf \cdot \Relrec \cdot \Opp \leq \frac{1}{\gamma}\left(\frac{\Delta}{\lambda} + \frac{\kappa(\Delta+1)}{\lambda^2} + \frac{\kappa^2}{\lambda^3}\right)
\label{eq:borne_forte_classique}
\end{equation}
dont le terme dominant $\Delta/(\gamma\lambda)$ est en $O(\lambda^{-1})$ et les
termes correctifs en $O(\lambda^{-2})$, sous l'Hypoth\`ese~\ref{ass:regime}
($I(M;\Sigma|C) \leq \kappa/\lambda$). Le Trilemme CRO post-quantique
(Contribution C3, section~\ref{sec:cro_pq}) \'etablit la borne correspondante
\begin{equation}
\Conf^q \cdot \Relrec^q \cdot \Opp^q \leq \frac{1}{\gamma_q}\left(\frac{\Delta_q}{\lambda} + \frac{\kappa(\Delta_q+1)}{\lambda^2} + \frac{\kappa^2}{\lambda^3}\right)
\label{eq:borne_forte_quantique}
\end{equation}
o\`u $\Delta_q$ repr\'esente la capacit\'e d'extraction institutionnelle
adapt\'ee au quantum et $\eta_q$ le coefficient de perte d'interpr\'etabilit\'e
quantique.

\begin{remark}[Statut d'une borne forte en $O(\lambda^{-2})$]
\label{rem:borne_forte_statut}
La borne~\eqref{eq:borne_forte_classique}, v\'erifi\'ee dans
l'annexe~\ref{chap:annexe_lean4}, a un terme dominant $\Delta/(\gamma\lambda)$ en
$O(\lambda^{-1})$~; seuls ses termes correctifs, $\kappa(\Delta+1)/\lambda^2$ et
$\kappa^2/\lambda^3$, sont d'ordre $\lambda^{-2}$ et $\lambda^{-3}$. Une borne dont
le terme dominant serait en $O(\lambda^{-2})$ reste une question ouverte.
La forme raffin\'ee~\eqref{eq:trilemme_cro_raffine} montre d'ailleurs qu'en
r\'egime \og rare\fg{} la borne est gouvern\'ee par Fano, en $1/\lambda$, ind\'ependamment
de $\Delta$. Les autres probl\`emes ouverts concernent l'extension de la
formalisation \`a SOTP et au cadre post-quantique ainsi que la d\'erivation dans
Mathlib des in\'egalit\'es de Fano et Pinsker, aujourd'hui pos\'ees comme
propri\'et\'es des instances.
\end{remark}

La formalisation de l'optimisation analytique sur l'intervalle complet
$[0, \lambda]$ reste un probl\`eme ouvert, car elle n\'ecessite des outils
d'analyse r\'eelle (d\'eriv\'ees de fonctions avec racine carr\'ee) non encore
disponibles dans Mathlib. Elle permettrait d'\'etudier le comportement de la
borne hors du r\'egime de l'Hypoth\`ese~\ref{ass:regime}.

\subsubsection{G\'en\'eralisation \`a la r\'egion r\'ealisable compl\`ete}
\label{subsubsec:region_realisable}

La borne sur le produit scalaire ne r\'epond pas \`a la question
de savoir quelle est la r\'egion r\'ealisable compl\`ete
$\mathcal{R} \subset [0,1]^3$ des triplets $(\Conf, \Rel, \Opp)$
atteignables~\cite{MINKA2026CLO}. Deux syst\`emes de produit CRO
identique~($= 0{,}4$) peuvent pr\'esenter des profils radicalement
diff\'erents~; seule la caract\'erisation de $\mathcal{R}$ permet de
d\'eterminer quels profils sont r\'ealisables et quels sont les arbitrages
Pareto-optimaux~\cite{Boyd2004}.

Deux conjectures guident l'exploration. La premi\`ere, dite de \emph{convexit\'e}, pose que $\mathcal{R}$
est convexe dans $[0,1]^3$, car la composition probabiliste de deux protocoles
g\'en\`ere un troisi\`eme avec m\'etriques interpol\'ees~\cite{MINKA2026CLO}.
La seconde, dite de \emph{fronti\`ere poly\'edrale}, pose que la fronti\`ere de Pareto de $\mathcal{R}$
est d\'efinie par un syst\`eme fini d'in\'egalit\'es lin\'eaires issues de
Fano~\cite{Fano61}, Pinsker~\cite{Pinsker1964} et du principe de traitement
des donn\'ees~\cite{CoverThomas06}. Des techniques issues de la programmation
convexe~\cite{Boyd2004,Nesterov2018} et de la th\'eorie taux-distortion
\'etendue~\cite{Berger71,Gray2011} semblent prometteuses pour attaquer ces
questions.

\subsubsection{G\'en\'eralisation au chiffrement homomorphe (FHE)}
\label{subsubsec:generalisation_fhe}

Le chiffrement totalement homomorphe~\cite{Gentry09,Vaikuntanathan2012,Brakerski2014}
maximise $\Conf$, les donn\'ees restent chiffr\'ees tout au long du calcul,
mais d\'egrade conjecturalement $\Opp$, car l'interpr\`ete institutionnel ne peut
extraire d'information verdictuelle d'un chiffr\'e homomorphe sans cl\'e de
d\'echiffrement~\cite{Gentry2010,Halevi2012}. Cette tension n'est pas encore
formalis\'ee dans le cadre CLO. La r\'esolution pourrait passer par des
constructions autorisant un d\'ecryptage partiel
contr\^ol\'e~\cite{Brakerski2013,Boneh2013} vers la couche Clay sans
compromettre la couche Gold~; il s'agit d'une conjecture.

\subsubsection{G\'en\'eralisation au calcul multi-parties s\'ecuris\'e (MPC)}
\label{subsubsec:generalisation_mpc}

Dans un protocole MPC \`a $n$
parties~\cite{Yao82,GMW87,Goldreich1987,Evans2018,Garay15}, l'opposabilit\'e
agr\'eg\'ee
\begin{equation}
  \Opp_\text{MPC} = \frac{I(V_1,\ldots,V_n;\,\Sigma \mid C)}{H(V_1,\ldots,V_n \mid C)}
\end{equation}
peut exhiber des comportements radicalement diff\'erents de l'opposabilit\'e
centralis\'ee~\cite{MINKA2026CLO}. Des effets d'amplification~(chaque partie
apportant une vue partielle enrichissant l'information de verdict) ou de
cancellation~(vues partielles contradictoires r\'eduisant
$I(V_1,\ldots,V_n;\,\Sigma \mid C)$) sont th\'eoriquement envisageables, sous
r\'eserve de formalisation. La structure du Trilemme CRO se pr\'eserve-t-elle
dans ce cadre~? Des techniques issues du mod\`ele
UC~\cite{Canetti01,Canetti2000} et de la th\'eorie des jeux
coop\'eratifs~\cite{Peleg2007,Owen2013} semblent indiqu\'ees.

\subsubsection{Cadre quantique quantitatif}
\label{subsubsec:cadre_quantique_quantitatif}

La migration vers des constructions quantiques, commitments bas\'es sur les
r\'eseaux~\cite{Ajtai96,Regev05,Peikert2016},
STARKs~\cite{STARKs18,BenSasson2018}, pr\'eserve-t-elle les propri\'et\'es
de confidentialit\'e et de fiabilit\'e de v\'erification \'etablies \`a la
section~\ref{sec:sotp_construction}, sachant que les engagements de Pedersen
reposent sur le logarithme discret et ne sont pas post-quantiques~? Quelles bornes CLO quantitatives
s'obtiennent pour des adversaires Grover~\cite{Grover96,Zhandry2012}~?
L'opposabilit\'e naturelle pour des protocoles \`a $r$ rounds
\begin{equation}
  \Opp_r = \frac{I(V;\,\Sigma_1,\ldots,\Sigma_r \mid C)}{H(V \mid C)}
\end{equation}
satisfait-elle une variante du Trilemme CRO~? Ces questions prolongent l'\'etape B
de la progression, par l'\'etablissement de bornes g\'en\'eralis\'ees.

L'extension quantique (Contribution C3, section~\ref{sec:cro_pq}) a ouvert
cette voie en introduisant les param\`etres $\Delta_q$ et $\eta_q$. Le
programme de recherche vise maintenant \`a calibrer ces param\`etres pour les
signatures NIST PQC (ML-DSA~\cite{Dilithium2022}, SLH-DSA~\cite{SPHINCS2022}),
pour lesquelles $\eta_q = 0$, puis \`a d\'evelopper une th\'eorie quantitative des
pertes d'interpr\'etabilit\'e pour le sc\'enario d'un support quantique~\cite{NielsenChuang2010,Preskill1998}.

Trois pistes th\'eoriques prolongent cette extension. \textbf{(D) Contraction.}
Pour un canal $\Lambda$ de coefficient de contraction $\eta_\Lambda$
(in\'egalit\'e de traitement des donn\'ees forte), on aurait
$\Opp^q \leq \eta_\Lambda\,\Opp$, ce qui remplacerait le terme additif
$\eta_q H(C)$ par un facteur multiplicatif. \textbf{(E) Granularit\'e et
capacit\'e.} Lier $\alpha$ et $\Delta$ (un interpr\`ete plus capable distingue
des verdicts plus fins) ferait d\'ependre $\gamma$ de $\Delta$ et modifierait le
r\'egime de la forme raffin\'ee. \textbf{(B2) Contexte priv\'e.} Pour SOTP, un
mod\`ele o\`u une partie $\tilde C$ du contexte est transmise \`a l'institution
seule (et non publique) permettrait de donner un statut prouv\'e aux
minorations d'opposabilit\'e, aujourd'hui heuristiques.

\subsubsection{Enseignements m\'ethodologiques de la formalisation Lean4}
\label{subsubsec:enseignements_lean4}

La formalisation Lean4 du noyau alg\'ebrique du Trilemme CRO
(section~\ref{sec:lean4}) a r\'ev\'el\'e trois le\c{c}ons structurelles pour
le d\'eveloppement du sous-domaine.

\paragraph{Propri\'et\'es informationnelles comme conditions v\'erifiables.}
La formalisation int\`egre Fano et Pinsker comme champs explicites de la
structure \texttt{ProtocolInstance} tandis qu'elle traite l'extraction born\'ee et le r\'egime de
confidentialit\'e comme hypoth\`eses du th\'eor\`eme. Ce choix transforme des
propri\'et\'es profondes en conditions v\'erifiables \`a chaque instance,
clarifiant le statut \'epist\'emique du th\'eor\`eme.

\paragraph{Coh\'erence et port\'ee du th\'eor\`eme.}
La formalisation produit deux t\'emoins qui d\'emontrent la non-vacuit\'e et
la n\'ecessit\'e des hypoth\`eses. Le \emph{t\'emoin de coh\'erence} atteint
exactement la borne th\'eorique pour une instance du mod\`ele abstrait,
confirmant que les hypoth\`eses formelles sont compatibles (ce point n'est
pas r\'ealisable par un protocole, Remarque~\ref{rem:portee_temoin}). Le \emph{t\'emoin de port\'ee} v\'erifie
les Hypoth\`eses 1, 2$'$ et 3 mais s'\'ecarte du r\'egime de confidentialit\'e,
atteignant $1/4$ constant, ce qui \'etablit que ce r\'egime est n\'ecessaire
\`a la borne. Le cas limite complet est document\'e \`a la
sous-section~\ref{subsec:lean4_erreurs}.

\paragraph{Couverture du corpus et scalabilit\'e.}
La formalisation couvre le noyau alg\'ebrique du th\'eor\`eme.
L'extension aux constructions SOTP compl\`etes et aux r\'esultats
post-quantiques constitue une direction de recherche future, estim\'ee \`a
un investissement substantiel mais r\'ealiste pour une \'equipe d\'edi\'ee.

\subsubsection{Extension de la formalisation Lean4~: SOTP et PQ}
\label{subsubsec:extension_lean4_sotp_pq}

La formalisation Lean4 actuelle (annexe~\ref{chap:annexe_lean4}) couvre
int\'egralement le noyau alg\'ebrique du Trilemme CRO
classique, y compris les t\'emoins de coh\'erence et de port\'ee. Deux
extensions restent \`a d\'evelopper.

\begin{itemize}
    \item \textbf{Extension SOTP.} La formalisation des constructions
          SOTP-Low, SOTP-Base et SOTP-High n\'ecessiterait la
          formalisation des r\'eductions cryptographiques (PRG, DLog,
          q-SDH) et des algorithmes de compression
          Johnson-Lindenstrauss. L'effort est estim\'e \`a 1000--1500 lignes
          additionnelles.
    \item \textbf{Extension post-quantique.} La formalisation du
          Trilemme CRO PQ n\'ecessiterait l'int\'egration des
          in\'egalit\'es quantiques (Pinsker quantique, Gentle
          Measurement, Fano quantique) dans une biblioth\`eque de
          th\'eorie de l'information quantique. L'effort est estim\'e \`a
          1000--1500 lignes additionnelles.
\end{itemize}

Cette extension, estim\'ee \`a 2000--3000 lignes additionnelles au total,
constituerait la premi\`ere certification m\'ecanique compl\`ete d'un
cadre cryptographique institutionnel.
\end{sloppypar}

% -----------------------------------------------------------------------
\subsection{Axe IV. Cryptographie institutionnelle et intelligence artificielle}
\label{subsec:axe4_ia}
% -----------------------------------------------------------------------
\begin{sloppypar}

\subsubsection{Interpr\`etes institutionnels artificiels}
\label{subsubsec:interpretes_artificiels}

Le cadre CLO mod\'elise l'interpr\`ete institutionnel par sa capacit\'e
cognitive $\Delta$ et un espace de verdicts $\mathcal{V}$~\cite{MINKA2026CLO,Simon82}.
Un mod\`ele de langage de grande taille~(LLM) utilis\'e comme interpr\`ete
pr\'esente des propri\'et\'es qualitativement diff\'erentes, puisque $\Delta_\text{LLM}$
peut d\'epasser $H(\mathcal{C})$ gr\^ace \`a des connaissances
hors-contexte~\cite{Brown2020,OpenAI2023}, mais est \'egalement stochastique
et potentiellement adversarial~\cite{Carlini2021,Zou2023}. Si $\Delta \gtrsim
H(\mathcal{C})$, la borne CRO (forme raffin\'ee~\eqref{eq:trilemme_cro_raffine}) peut-elle
\^etre approch\'ee arbitrairement~? Autrement dit, des interpr\`etes IA
rel\`event-ils fondamentalement la fronti\`ere de Pareto dans l'espace
CLO~? Cette question est ouverte~\cite{Bommasani2021}. Un adversaire IA, LLM
finement ajust\'e pour l'extraction d'information institutionnelle,
constitue-t-il une menace qualitativement nouvelle pour la couche Clay, ou
reste-t-il born\'e par les limites issues de la th\'eorie de
l'information~\cite{CoverThomas06}~? Les deux questions rel\`event de l'\'etape F
du programme, car elles requi\`erent d'abord une d\'efinition formelle de
$\Delta_\text{LLM}$ avant toute borne~\cite{Weidinger2023}.

\subsubsection{Apprentissage automatique comme outil d'estimation CLO}
\label{subsubsec:apprentissage_estimation}

Les m\'etriques CLO sont d\'efinies au sens de la th\'eorie de
l'information~\cite{Shannon48,CoverThomas06}; leur estimation sur donn\'ees
r\'eelles est computationnellement difficile~\cite{Paninski2003}. Les
techniques d'estimation d'information mutuelle par apprentissage~(MINE~\cite{Belghazi18},
m\'ethodes par diff\'erence de log-vraisemblance~\cite{Poole2019,Oord2018})
ouvrent une voie vers des \'evaluations empiriques automatis\'ees des
propri\'et\'es CLO. L'estimation de $H(\mathcal{C})$ pour des contextes
institutionnels r\'eels via des mod\`eles de compression
neuronaux~\cite{Hinton2007,Kingma2013}, l'estimation de $\Delta$ par des tests
de compr\'ehension automatis\'es sur corpus de d\'ecisions institutionnelles,
ces directions permettraient de remplacer les estimations analytiques
actuelles~(ECDSA, Helios, zk-SNARKs) par des mesures directes. Leurs
r\'esultats resteraient soumis \`a calibration
exp\'erimentale~\cite{Campbell1976,Saltelli2008}; ils constitueraient
n\'eanmoins la premi\`ere validation empirique directe du cadre.

\subsubsection{Audit cryptographique de l'intelligence artificielle}
\label{subsubsec:audit_ia}

La direction inverse est \'egalement possible. Des protocoles SOTP permettant
de prouver des propri\'et\'es de mod\`eles IA,
robustesse~\cite{Goodfellow2015,Madry2018},
\'equit\'e~\cite{Barocas2019,Mitchell2021}, conformit\'e
r\'eglementaire~\cite{EU2024AI,EU2021AI}, \`a des r\'egulateurs
institutionnels sans r\'ev\'eler les poids du mod\`ele constituent une
application naturelle du paradigme CLO. L'audit cryptographique d'un mod\`ele
de d\'ecision consisterait ainsi \`a prouver sa conformit\'e via un sch\'ema SOTP-High sans exposer
son architecture. Ce programme croise la v\'erification du
calcul~(\textit{verifiable computation})~\cite{Gennaro2010,Walfish2015} et la
gouvernance de l'IA~\cite{Floridi2018,Jobin2019}, territoire d'application
neuf pour lequel le cadre CLO-CRO-SOTP est th\'eoriquement arm\'e.
\end{sloppypar}

% -----------------------------------------------------------------------
\subsection{Programme de recherche \`a horizon structur\'e}
\label{subsec:programme_recherche}
% -----------------------------------------------------------------------
\begin{sloppypar}

\begin{boiteArbitre}
L'arbitre quitte l'exploration. Quatre axes ont \'et\'e ouverts, qui portent sur les fondations
computationnelles, la voie protocolaire, les extensions th\'eoriques et l'intelligence
artificielle. Chacun ouvre des territoires prometteurs, mais en partie encore sp\'eculatifs.
Il faut d\'esormais planifier, en distinguant ce qui peut \^etre
fait en un an, en trois ans ou en dix ans selon que cela d\'epend d'une preuve, d'un
financement ou d'une communaut\'e. L'arbitre note que le propos rel\`eve d\'esormais moins d'un
chercheur isol\'e que d'un programme qui se donne des
\'ech\'eances.
\end{boiteArbitre}

\subsubsection{Horizon court terme (1--3 ans)~: consolidation des fondations}
\label{subsubsec:horizon_court_terme}

Cinq t\^aches conditionnent les horizons ult\'erieurs. D'abord,
\'ecrire l'argument qui exclut l'inversion successive par extraction de
racines, puis soumettre AIIP \`a la cryptanalyse~; la robustesse de cette
hypoth\`ese est pr\'ealable \`a tout usage applicatif du CASH
Toolkit~\cite{MINKA2025AIIP}. Ensuite,
\'etendre la formalisation Lean4 \`a SOTP et au cadre post-quantique
(cf.~sous-section~\ref{subsubsec:extension_lean4_sotp_pq}), avec pour objectif
d'\'elargir la couverture de l'annexe~\ref{chap:annexe_lean4} sans modifier la
borne. Parall\`element,
exhiber la r\'eduction UC entre SOTP-High et ZK-NR~v3, ce qui fermerait
formellement la boucle entre la th\'eorie abstraite de
SOTP~(section~\ref{sec:sotp_construction}) et son instanciation
composable~\cite{MINKA2025ZKNR_UC,Canetti01}. Puis, prouver un
th\'eor\`eme de composition garantissant que l'assemblage
$(CEE \;\|\; AOW \;\|\; SH)$ en Gold-Iron-Clay pr\'eserve les bornes CRO
\`a un facteur constant. Enfin, ex\'ecuter le
programme exp\'erimental sp\'ecifi\'e au
chapitre~\ref{chap:experimentation}.

Les six protocoles y sont d\'ecrits jusqu'au niveau op\'erationnel (stratification des \'echantillons, estimateurs, crit\`eres de fiabilit\'e
inter-\'evaluateurs, seuils de falsification) mais aucun n'a \'et\'e
d\'eploy\'e (sous-section~\ref{subsubsec:statut_deploiement}). Tant que
$\Delta$, $H(\mathcal{C})$, $\Delta_q$ et $\eta_q$ demeurent estim\'es
par triangulation documentaire, sans mesure aupr\`es d'acteurs
institutionnels r\'eels, les bornes CRO conservent le statut de
r\'esultats conditionnels, puisqu'elles sont math\'ematiquement \'etablies sous
hypoth\`eses, sans calibration empirique. Le d\'eploiement suppose trois pr\'ealables non
techniques, \`a savoir l'accord institutionnel des juridictions et r\'egulateurs
approch\'es, l'approbation d'un comit\'e d'\'ethique de la recherche et un
financement sur un horizon de vingt-six mois. L'audit post-\'electoral
Helios de la section~\ref{subsec:protocole_opp}, le plus l\'eger en
ressources et le seul enti\`erement internalisable \`a l'ENSP, constitue
le point d'entr\'ee naturel, puisqu'il valide $\Opp$ sur un syst\`eme
r\'eellement d\'eploy\'e et fournit le premier point de calibration du
cadre.

\subsubsection{Horizon moyen terme (3--5 ans)~: \'elargissement du domaine}
\label{subsubsec:horizon_moyen_terme}

La th\'eorie CLO computationnelle, d\'efinitions de
$(\Conf_\text{comp}, \Rel_\text{comp}, \Opp_\text{comp})$ via adversaires
PPT~\cite{Goldreich01,Bellare97} et \'etablissement du trilemme dans ce cadre
renforc\'e, est la priorit\'e th\'eorique centrale de cet horizon. En
parall\`ele, les variantes du Trilemme CRO pour FHE et
MPC~\cite{Gentry09,Evans2018} doivent \^etre d\'emontr\'ees ou r\'efut\'ees~;
les bornes d'opposabilit\'e atteignables dans ces extensions restent
inconnues. Des outils d'estimation de $H(\mathcal{C})$ et $\Delta$ par
apprentissage automatique peuvent \^etre d\'evelopp\'es et calibr\'es sur des
preuves existantes dont les m\'etriques CLO ont \'et\'e estim\'ees
analytiquement ($0{,}437$ pour Helios, $0{,}572$ pour zk-SNARKs et $0{,}273$ pour ECDSA), sous
r\'eserve que les estimations analytiques constituent une r\'ef\'erence
raisonnablement proche des mesures directes. La contribution \`a la
standardisation, int\'egration de l'opposabilit\'e dans les standards NIST
post-quantiques, cadrage d'un futur standard ISO/IEC, peut d\'ebuter sur cet
horizon.

\subsubsection{Horizon long terme (5--10 ans)~: sous-domaine \'etabli}
\label{subsubsec:horizon_long_terme}

La vision d\'ecennale est celle d'un sous-domaine reconnu dans les
conf\'erences de cryptologie~(CRYPTO, EUROCRYPT, ASIACRYPT), avec sa propre
journ\'ee de travail sp\'ecialis\'ee, des programmes doctoraux d\'edi\'es dans
les universit\'es de recherche ainsi que des venues
propres~(\textit{International Symposium on Institutional Cryptography},
\textit{Journal of Institutional Cryptography}). Scientifiquement, l'horizon
d\'ecennal vise la caract\'erisation compl\`ete de $\mathcal{R} \subset [0,1]^3$,
fronti\`ere de Pareto exacte comprise, une th\'eorie unifi\'ee couvrant
protocoles non interactifs, interactifs et multi-parties sous un cadre
CLO-CRO \'etendu. Si AIIP r\'esiste \`a la cryptanalyse, cet horizon inclut aussi des primitives
CASH soumises \`a standardisation.

\begin{table}[htbp]
\centering
\caption{Synth\`ese des perspectives par contribution et horizon}
\label{tab:perspectives_horizons}
\small
\begin{tabular}{@{}p{2.5cm}p{3cm}p{3cm}p{3cm}@{}}
\toprule
\textbf{Contribution} & \textbf{Court terme} & \textbf{Moyen terme} & \textbf{Long terme} \\
\midrule
C1 (CLO) & D\'eploiement du protocole (chap.~\ref{chap:experimentation}, non ex\'ecut\'e) & Outils ML d'estimation & Standards ISO/IEC \\
C2 (CRO classique) & Extension Lean4 \`a SOTP et PQ & Extensions FHE/MPC & Caract\'erisation $\mathcal{R}$ \\
C3 (CRO post-quantique) & Mesure de $\Delta_q$ (ML-DSA, SLH-DSA) & Th\'eorie $\eta_q$ (options D, E) & PQC certifi\'e CLO \\
C4 (SOTP) & Mod\`ele \`a contexte priv\'e (B2), R\'eduction SOTP$\to$ZK-NR & Preuve ou r\'efutation de la conjecture d'optimalit\'e & Th\'eorie unifi\'ee, impl\'ementations de r\'ef\'erence \\
C5 (Lean4) & Extension couverture SOTP et PQ & Certification UC & V\'erificateur automatique, int\'egration Mathlib et Coq/Isabelle \\
C6 (Protocoles) & D\'eploiement Helios audit post-\'electoral & Validation multi-institutionnelle & Standards empiriques \\
\bottomrule
\end{tabular}
\end{table}

\subsubsection{Financements envisag\'es}
\label{subsubsec:financements_envisages}

\noindent\textbf{Agences nationales.} MINRESI (Cameroun), ANR (France).

\noindent\textbf{Agences internationales.} NSF SaTC (\'Etats-Unis), IDRC
(Canada), Horizon Europe (EU).

\noindent\textbf{Partenariats institutionnels.} LIMSI (Cameroun), NIST (PQC
validation), IBM Quantum (acc\`es simulateur).

\subsubsection{Suivi longitudinal}
\label{subsubsec:suivi_longitudinal}

Les travaux ant\'erieurs~\cite{MINKA2026CRO} et la
section~\ref{sec:cro_pq} invitent \`a
mesurer l'\'evolution des param\`etres \`a mesure que les institutions
acqui\`erent de l'exp\'erience, selon les lois
\begin{align}
\Delta(t) &= \Delta_\infty(1 - e^{-t/\tau})
\label{eq:suivi_delta} \\
\Delta_q(t) &= \Delta_{q,\infty}(1 - e^{-t/\tau_q})
\label{eq:suivi_delta_q} \\
\eta_q(t) &= \eta_{q,0} \cdot e^{-t/\tau_{\mathrm{hw}}}
\label{eq:suivi_eta_q}
\end{align}
o\`u $\tau$ est la constante de temps d'apprentissage institutionnel et
$\tau_{\mathrm{hw}}$ refl\`ete l'am\'elioration du mat\'eriel quantique. Ce
programme longitudinal repr\'esente une direction de recherche \`a 3--5 ans.
\end{sloppypar}

\begin{boiteArbitre}
L'arbitre retient que le voyage ne s'arr\^ete pas ici. Les bornes existent
sous hypoth\`eses, tandis que les constructions ont des propri\'et\'es prouv\'ees et une
opposabilit\'e encore heuristique. Une borne de terme dominant
$O(\lambda^{-2})$ reste ouverte, de m\^eme que les extensions au calcul homomorphe
et multi-parties~; les param\`etres quantiques
$\Delta_q$ et $\eta_q$ restent quant \`a eux \`a calibrer sur des acteurs r\'eels. Ce qui a
\'et\'e construit est prouv\'e l\`a o\`u il est dit prouv\'e et conjectur\'e l\`a o\`u il est dit conjectur\'e.
Ce qui reste \`a construire est plus vaste encore, mais une carte existe d\'esormais pour s'y orienter.
Un sous-domaine est en marche avec ses questions ouvertes, ses hypoth\`eses
\`a valider et sa communaut\'e qui commence \`a se constituer autour d'elles.
\end{boiteArbitre}